% Physical unit 033; accepted Session 03 after fresh 400-dpi authority replay.
\GordanUnitStart{033}{printed page 28}
\GordanTopFolio{28}
\begingroup\small
\sect{5}

For computing $K$ we use the formula
\[
 \sqrt{2K\over\pi}=1+2q \qquad \text{fund. p. 184}
\]
or
\[
 K={\pi\over2}\{1+2q\}^2.
\]
For computing $E$, we use the formulae
\[
\begin{array}{rcll}
 {2K\over\pi}{2K\over\pi}-{2K\over\pi}\cdot {2E\over\pi}&=&8q+16q^2+32q^3\ldots, &135\vspace{0.2em}\\
 {2K\over\pi}\cdot{2K\over\pi}&=&1+8q+24q^2+32q^3\ldots, &105
\end{array}
\]
By subtracting one obtains:
\[
 {2K\over\pi}\cdot {2E\over\pi}=1+8q^2,
\]
and by dividing by the preceding equation:
\[
 {2E\over\pi}={1+8q^2\over (1+2q)^2}.
\]
This is the other method of computing $\Sigma(l_0)$.

For there is:
\[
 \Sigma(l_0)={b\over x'}E \qquad \text{Cap. VI \S\ 3}
\]
whence
\[
 {2\over\pi}\Sigma(l_0)=b{1+8q^2\over \{1-2q\}^2}.
\]

\sect{6}

It remains for $v$ to be computed; since this is not explicitly available, it is useful to compute the trigonometric functions related to it of the fraction
\[
 {iv\pi\over2K}.
\]
Accordingly we wish to introduce a new constant $A$, to be defined by the equation
\[
 \sin A\cos {iv\pi\over2K}=1
\]
or:
\[
 \sin {iv\pi\over2K}=i\cot A
 \qquad
 \tg {iv\pi\over2K}=i\cos A,
\]
from which it is understood that only $A$ has to be computed.
\endgroup
